\documentclass[12pt]{article}
\usepackage[utf8]{inputenc}
\usepackage{amsmath, amssymb} % Math packages
\usepackage{geometry}
\usepackage{hyperref} % Clickable links
\usepackage{graphicx} % Images

\geometry{a4paper, margin=1in}

\title{reti di calcolatori appunti}
\author{Matteo}
\date{\today}

\begin{document}


\maketitle
\section{secondo semestre appunti}
\section{sicurezza  di rete}
piu grande rischio di sicurezza è avere attaccante ma non saperlo

la crittografia e la piu grande arma che abbiamo.
volia garantire sia la cifrazione del documento che l'autenticazione dell'interlocutore che l'integrita del messaggio

firewall sta subto dopo il router e filtra tutto cio che arriva

si possono usare socket cifrati a livello trasporto, se no si puo cifrare anche a livello rete.

cufratura e la prima risorsa di sicurezza che viene usata

chiavi di cifratura usate per risolvere le cifrature. ce una chiave giusta, piu un numero esorbitante di chiavi esca.

esiste l'encryption algorithm(sender) e il decryption algorithm(receiver); la chiave serve al receiver.

la chiave e condivisa dal sender insieme alle esche.

DES= standard di cifrazione, era a 56 bit, fu fatta una challenge e fu hakerato in meno di un giorno, cosi il numero di bit fu aumentato di 3 volte.




\end{document}